% Partie : histoire
% Chapitre : biographies
% Section : Georg Cantor

\documentclass[../../../include/open-logic-section]{subfiles}

\begin{document}

\olfileid[fr]{his}{bio}{can}

\olsection{Georg Cantor}

\olphoto{cantor-georg}{Georg Cantor}

Une ancienne biographie
de Georg Cantor
(\textsc{gué}-org
\textsc{kan}-tor)
affirmait qu'il
était né et
avait été trouvé
à bord d'un
navire faisant
route vers
Saint-Pétersbourg,
en Russie, et
que ses parents
étaient inconnus.
Ce récit est
faux ; Cantor
est toutefois
né à
Saint-Pétersbourg
en 1845.

Cantor obtint
son doctorat
en mathématiques
à l'université
de Berlin
en 1867.
Il est connu
pour ses travaux
sur la théorie
des ensembles
et considéré
comme le fondateur
de ce domaine
de recherche
distinct. Il
fut le premier
à démontrer
qu'il existe
des ensembles
infinis de
cardinaux différents.
Ses théories,
en particulier
sa théorie
des infinis,
suscitèrent
de nombreux
débats parmi
les mathématiciens
de son époque
et rendirent
ses travaux
controversés.

Les convictions
religieuses de
Cantor et ses
travaux mathématiques
étaient intimement
liés. Il affirma
même que la
théorie des
nombres transfinis
lui avait été
communiquée
directement
par Dieu.
Vers la fin
de sa vie,
Cantor souffrit
de troubles
psychiques.
Il fut hospitalisé
à partir de 1894,
puis de plus
en plus souvent
dans ses dernières
années. Ses travaux suscitèrent de vives critiques,
notamment de la part du mathématicien Leopold Kronecker,
avec lequel il se brouilla. Pendant ses épisodes dépressifs,
Cantor se détourna parfois des mathématiques et se tourna
vers la philosophie et la littérature.
Il publia même
une théorie
selon laquelle
Francis Bacon
était l'auteur
des pièces
de Shakespeare.

Cantor mourut
le 6 janvier 1918
dans un sanatorium
à Halle.

\begin{reading}
Pour des biographies
complètes de Cantor,
voir \citet{Dauben1990}
et
\citet{Grattan-Guinness1971}.
Ses idées radicales
sont aussi présentées
dans l'émission de
BBC Radio~4
\emph{A Brief History
of Mathematics}
\citep{Sautoy2014}.
Pour découvrir
les théories
de Cantor sous
forme de rap,
voir \citet{Rose2012}.
\end{reading}

\end{document}
